%%%BLOCK id=001-01 ch=90 sec=微分方程·分离变量 kind=derivation alt=none cont=none y=0.06-0.13
\[
v=kx,\quad \frac{\dd x}{\dd t}=kx,\quad \int\frac{\dd x}{kx}=\int\dd t,\quad \ln x=kt+kc,\quad x=\ee^{kt+kc}
\]
%%%END
%%%BLOCK id=083-07 ch=90 sec=定义式与决定式 kind=summary alt=01,02,03,09,10,11 cont=none y=0.11-0.23
% 原稿左侧有大括号：{定义式 比值 (无关)；决定式 (相关)}，“无关”“相关”均被圈起；公式按原稿行序排列
\[
\left\{\begin{aligned}
&\text{定义式\ \ 比值}\\[8pt]
&\text{决定式}
\end{aligned}\right.
\]

$v=\dfrac{\Delta x}{\Delta t}$\quad $k=\dfrac{F}{x}$\quad $E=\dfrac{F}{q}$\quad $\unclear{R=\dfrac{U}{I}}$ % CHECK: 末式笔迹形似“k=v/L”，按R=U/I转录(本页“R”常写得像“k”)

$I=\dfrac{q}{t}$\quad $C=\dfrac{Q}{U}$\quad $\Phi=B_{\perp}S$

\fbox{无关}\quad $p=mv$\quad $I=Ft$

\fbox{相关}\quad $a=\dfrac{F}{m}$\quad $E=\dfrac{kQ}{r^2}$

$I=nesv$\quad $R=\dfrac{\rho L}{S}$\quad $C=\dfrac{\varepsilon S}{4\pi kd}$
%%%END
%%%BLOCK id=083-08 ch=90 sec=函数分析·变化率与面积 kind=formula alt=01,06,09,10,12 cont=none y=0.12-0.40
\textbf{函数分析}

% 原稿左侧红笔大括号括住下列各式
\[
\left\{\begin{aligned}
a&=\frac{\Delta v}{\Delta t}\\[6pt]
E_{\text{\unclear{感}}}&=-n\frac{\Delta\Phi}{\Delta t}\qquad E_{\text{场}}=-\red{\frac{\Delta\varphi}{\Delta x}}\\[6pt]
i&=\frac{\Delta q}{\Delta t}\\[6pt]
F&=\frac{\Delta p}{\Delta t}\\[6pt]
P&=\frac{\Delta W}{\Delta t}\\[6pt]
F_{\text{合}}&=\frac{\Delta E_k}{\Delta x}\\[6pt]
F_{\text{其他}}&=\frac{\Delta E_{\text{机}}}{\Delta x}\\[6pt]
R&=\frac{U_1}{I_1}=\frac{U_2}{I_2}\neq\frac{U_2-U_1}{I_2-I_1}\\[6pt]
P&=U_1I_1\qquad\text{面积}\ S=\int xy
\end{aligned}\right.
\]
% 原稿E_场式中黑笔旧写的分式被红笔涂划，右侧红笔重写为Δφ/Δx；E_感下标笔迹似“源/感”

$\dfrac{\Delta v}{\Delta t}$\quad 匀圆变化\quad 方向变
%%%END
%%%BLOCK id=091-06 ch=90 sec=数列求和·多过程 kind=method alt=06 cont=none y=0.55-0.65
\textbf{数列类}
\[
\left\{\begin{aligned}
S_n&=\frac{(a_1+a_n)n}{2}\\
S_n&=\frac{a_1}{q-1}(q^{n}-1)\qquad \text{若}\ |q|<1\ \ \lim_{n\to\infty}S_n=\frac{a_1}{1-q}
\end{aligned}\right.
\]
\begin{quote}
详写过程 $|$ 易得过程2\\
综上全过程得出通项式\\
求前$n$项和、极限$n$无穷
\end{quote}
% CHECK: “过程”后的竖线也可能是数字 1（详写过程1）
%%%END
%%%BLOCK id=091-08 ch=90 sec=近似·二项式定理 kind=formula alt=none cont=none y=0.75-0.79
\textbf{二项式定理}\quad $(1+x)^{\alpha}\approx 1+\alpha x$\quad ($x$ 很小时)
%%%END
%%%BLOCK id=199-06 ch=90 sec=微积分·常用导数(tan x) kind=formula alt=none cont=none y=0.69-0.86
$(\tan x)'=\dfrac{1}{\cos^2 x}$

\notefig[0.55]{figures/p199_e.png}
% 图为 $y=\tan x$ 的导函数 $\frac{1}{\cos^2x}$ 的草图（竖直虚线为渐近线），箭头自公式指向中间一支；右侧另有一小段红笔画的 V 形曲线
%%%END
